\chapter{Experimental Campaign}
\label{chap:experimental_campaign}

\section{Overview}

The experimental campaign constitutes the foundation of the present work and aims to generate a comprehensive dataset for the investigation of forming limit curves (FLCs). Particular emphasis is placed on assessing the influence of specimen geometry, sheet thickness, punch geometry, material orientation, and loading path on the experimentally determined forming limits. The resulting dataset serves both for the evaluation of forming limit identification methodologies and for the validation of the numerical simulation framework developed within this project.

\section{Experimental Setup}

All experiments were conducted using the high-throughput forming test bench available at the Institute of Virtual Manufacturing (IVP). The system combines automated specimen handling with Digital Image Correlation (DIC) measurements, enabling the efficient acquisition of strain-field data for a large number of test configurations.

\begin{figure}[htbp]
    \centering
    \includegraphics[width=0.8\textwidth]{figures/test_bench.jpg}
    \caption{High-throughput forming test bench used for the experimental campaign.}
    \label{fig:test_bench}
\end{figure}

\section{Specimen Geometries and Test Configurations}

The experimental campaign is based on Nakazima, Marciniak, and Punch-in-Punch tests. By varying the specimen width, strain paths ranging from uniaxial tension to equi-biaxial stretching can be generated, enabling the construction of complete Forming Limit Diagrams.

Figure~\ref{fig:specimen_geometries} shows the investigated specimen geometries. Narrow specimens promote uniaxial deformation, whereas wider specimens progressively shift the strain state towards biaxial stretching.

\begin{figure}[htbp]
    \centering
    \includegraphics[width=0.9\textwidth]{figures/specimen_geometries.pdf}
    \caption{Investigated specimen geometries ranging from W20 to W200.}
    \label{fig:specimen_geometries}
\end{figure}

In addition to specimen geometry, the influence of sheet thickness, punch diameter, material orientation, and testing configuration was investigated. Besides conventional Nakazima and Marciniak tests, a Punch-in-Punch configuration was considered to introduce alternative loading paths and expand the range of achievable strain states.

\section{Experimental Matrix}

Table~\ref{tab:experimental_matrix} summarizes the main parameters investigated during the experimental campaign.

\begin{table}[htbp]
\centering
\caption{Overview of investigated experimental parameters.}
\label{tab:experimental_matrix}
\begin{tabular}{ll}
\hline
Parameter & Values investigated \\
\hline
Specimen geometry & W20--W200 \\
Sheet thickness & 1, 1.5, 2 ,3 [mm] \\
Punch diameter & 40, 60, 80, 100 [mm] \\
Material orientation & 0 (Rolling direction (RD)), 90 (Transverse direction (TD)) [deg] \\
Test configuration & Nakazima, Marciniak\\
\hline
\end{tabular}
\end{table}

\section{Data Acquisition}

Digital Image Correlation (DIC) was used to measure the full-field strain evolution throughout each test. The recorded strain histories constitute the primary input for the forming limit identification procedures presented in Chapter~\ref{chap:flc_identification}. To facilitate comparison with numerical simulations, both experimental and numerical datasets are organized around the same physical quantities, including major and minor strain histories, force--displacement responses, fracture location, and forming-limit points within the Forming Limit Diagram.